\documentclass[10pt,a4paper]{article}

\usepackage[T1]{fontenc}
\usepackage[utf8]{inputenc}
\usepackage{lmodern}
\usepackage{amsmath,amssymb,amsthm}
\usepackage{booktabs,array}
\usepackage{geometry}
\usepackage{xcolor}
\usepackage{enumitem}
\usepackage{microtype}
\usepackage{hyperref}
\usepackage{fancyhdr}

\geometry{margin=1in}
\hypersetup{
  colorlinks=true,
  linkcolor=blue!55!black,
  citecolor=blue!55!black,
  urlcolor=blue!55!black,
  pdftitle={Syndromes Solved Directly: A Zero-Signing-Query Forgery against SSD},
  pdfauthor={Martin Feussner},
  pdfsubject={Cryptanalytic audit report for IACR ePrint 2022/1698},
  pdfkeywords={digital signatures, syndrome decoding, affine completion, forgery, cryptanalysis}
}
\setlist{nosep,leftmargin=*}
\setlength{\parindent}{1em}
\setlength{\parskip}{0.25em}
\pagestyle{fancy}
\fancyhf{}
\fancyhead[L]{Syndromes Solved Directly}
\fancyhead[R]{14 October 2026}
\fancyfoot[C]{\thepage}
\renewcommand{\headrulewidth}{0.3pt}

\newtheorem{theorem}{Theorem}
\newtheorem{lemma}{Lemma}
\newcommand{\F}{\mathbb{F}}
\newcommand{\wt}{\mathsf{wt}}

\title{\bfseries Syndromes Solved Directly:\\
A Zero-Signing-Query Forgery against SSD}
\author{Martin Feussner\\
\small Selmer Center, University of Bergen\\
\small \texttt{martin.feussner@uib.no}}
\date{14 October 2026}

\begin{document}
\maketitle

\noindent\fcolorbox{black!50}{black!4}{\parbox{0.94\textwidth}{
This audit used the OpenAI model Daybreak Blue at the ultra reasoning level.
OpenAI Codex generated and selected this paper's title and intentionally chose
``Syndromes Solved Directly'' as wordplay that spells the scheme acronym SSD
and describes the direct public linear solve used by the forgery. OpenAI
Codex performed AI-assisted cryptanalysis, implementation, experimentation,
adversarial-review coordination, and technical writing. Primary reproduction,
a separately tasked hostile adversarial review, fresh-context independent
reproduction, the final exact-target attack-history recheck, and independent
Publication Reviewer review passed. No independent human verification was
performed or required.}}

\begin{abstract}
We give a public-key-only, zero-signing-query existential fresh-message forgery
against standard-mode SSD1760 from Abdelhaliem Babiker's \emph{Digital
Signature from Syndrome Decoding Problem}, IACR ePrint 2022/1698. The ordinary
verifier accepts $s\in\F_2^{1760}$ when
\[
  473\leq\wt(s)\leq513
  \quad\text{and}\quad
  Hs=Y\bar h.
\]
A valid $242\times1760$ public matrix $H$ has 242 pivots and 1518 free
coordinates. The attacker fixes 363 free coordinates to one and directly
solves the pivots for each public message target. In the paper's random-oracle
model, full-row-rank $Y$ makes the pivot vector uniform in $\F_2^{242}$.
Acceptance therefore occurs with exact probability
\[
  \Pr[110\leq\mathrm{Bin}(242,1/2)\leq150]
  =0.9303737782253141
\]
per distinct message, giving 1.0748368273 expected public hash queries and no
signing query. This is an existential message-search attack, not a universal
or fixed-preselected-message forgery.

The primary Python and hostile Magma reconstructions generated valid structured
keys, accepted honest signatures, and produced ordinary-verifier forgeries on
their first candidate messages. A fresh-context independent Python
reproduction then forged successfully for five of five independently generated
full-rank keys, using one or two candidate messages per key and zero signing
queries. The clean adversarial attack stage took 0.19 CPU seconds from before
row reduction; the independent attacks took 0.052984--0.053914 seconds per key
from the same boundary. Changed-message, changed-response, wrong-key, and
equation-valid weight-boundary controls behaved as required. The claim is
limited to valid full-rank standard SSD1760 keys in the random-oracle model.
LVK mode, other parameter rows, every fixed message, and an unconditional
SHAKE256 claim are outside scope.
\end{abstract}

\section{Introduction}

Signature from Syndrome Decoding (SSD) presents a binary hash-and-sign design
whose ordinary verifier checks a public linear equation and a Hamming-weight
interval~\cite{ssd}. The proposal argues that finding a response in this
interval requires solving a hard syndrome-decoding instance. For its first
advertised row, however, the public matrix has only 242 rows and the accepted
weight interval is centered broadly around half of those pivot coordinates
after a suitable public affine choice.

The attack uses standard Gaussian elimination. Put the public matrix in
systematic form, choose a fixed assignment to its free coordinates, and solve
the pivots for the public message syndrome. The resulting vector always
satisfies the equation. Its only remaining condition is weight. Fresh random-
oracle messages randomize the pivot assignment, and the advertised interval
accepts more than 93 percent of those assignments.

This distinction matters for the security claim. The attacker does not
control the random-oracle target and cannot promise success for every fixed
message. Instead, it chooses distinct fresh messages, computes a public
completion for each, and retains the first accepted pair. Existential
unforgeability permits the attacker to choose its forged message~\cite{gmr}, so this
constant-probability public search is sufficient. No message is submitted to
a signing oracle.

\paragraph{Contribution and provenance.}
We claim no novelty for Gaussian elimination, systematic form, affine linear
solving, or binomial probabilities. The result is their concrete application
to standard SSD1760's ordinary verifier at complete parameters. Primary
reproduction and a hostile adversarial review passed. Fresh-context
independent reproduction, final exact-target history review, and independent
publication review also passed. We make no priority claim and no technique
novelty claim.

\section{Target and exact scope}

\subsection{Pinned paper}

The target is Abdelhaliem Babiker's \emph{Digital Signature from Syndrome
Decoding Problem}, IACR ePrint 2022/1698~\cite{ssd}. The pinned current object
is a 458,702-byte, 14-page PDF with SHA-256 digest
\begin{center}
\small\ttfamily
619903383e40bed6320d1fc0a17c61553073d11545eb12302cafe6a84cfbd80e.
\end{center}
Its text extraction has SHA-256 digest
\begin{center}
\small\ttfamily
f2962618bf752d315aac0759e286491e6afffccf6d5764bd2c4c89f654c4a62b.
\end{center}
The official record observed on 14 October 2026 listed receipt on 7 December
2022 and approval on 10 December 2022. The version-list endpoint could not be
enumerated, so every claim is pinned to the bytes above.

No author implementation, concrete hash instantiation, message encoding, or
signature serialization was available. The executable evidence uses
independent paper-level reconstructions and an explicit SHAKE256 adapter. The
security theorem below uses the paper's stated random-oracle model and does
not depend on SHAKE256 behavior at every fixed message.

\subsection{Standard SSD1760 verifier}

The claimed row is standard SSD1760:
\[
 (n,k',d,\omega,\mu,k)=(1760,242,502,484,11,440),
 \qquad \gamma=4.
\]
The public matrices are
\[
 H\in\F_2^{242\times1760},
 \qquad
 Y\in\F_2^{242\times440}.
\]
The hash $h\in\F_2^{1760}$ is parsed into four 440-bit blocks
$h_1,\ldots,h_4$, and
\[
 \bar h=h_1+h_2+h_3+h_4\in\F_2^{440}.
\]
Algorithm~4 accepts a response $s\in\F_2^{1760}$ precisely when
\begin{equation}
 473\leq\wt(s)\leq513,
 \qquad Hs=Y\bar h.
 \label{eq:verify}
\end{equation}
There is no further distribution, support, proof, or encoding check in that
ordinary verifier.

The result is limited to valid public keys satisfying
$\operatorname{rank}(H)=\operatorname{rank}(Y)=242$, as the paper's key
generation requires. It covers standard mode only. LVK uses a secret subset
of public rows and is outside the attack model. SSD2744, SSD3728, malformed
keys, and rank-deficient Gen1 draws are also outside the claim.

\section{Public affine-completion forgery}

\subsection{Systematic form}

Public Gaussian elimination selects 242 pivot columns of $H$ and leaves
$1760-242=1518$ nonpivot coordinates. Write the corresponding decomposition
as
\[
 H_Ps_P+H_Fs_F=t,
 \qquad t=Y\bar h.
\]
The square pivot matrix $H_P$ is invertible. For any chosen free assignment
$s_F$ and target $t$, the unique pivot assignment is
\begin{equation}
 s_P=H_P^{-1}(t+H_Fs_F).
 \label{eq:completion}
\end{equation}
Thus every public target has an efficiently computable response satisfying
the verifier equation. The attack fixes an $s_F$ of weight 363.

\begin{lemma}[Uniform public target]
In the paper's random-oracle model, if $Y$ has rank 242, then $Y\bar h$ is
uniform in $\F_2^{242}$ for every fresh distinct message.
\end{lemma}

\begin{proof}
The xor map from four 440-bit blocks to one 440-bit vector is surjective, so it
maps a uniform 1760-bit random-oracle output to a uniform element of
$\F_2^{440}$. A rank-242 linear map $Y:\F_2^{440}\rightarrow\F_2^{242}$ is
surjective and maps that uniform input to a uniform 242-bit output. Distinct
random-oracle queries give independent outputs.
\end{proof}

\begin{theorem}[Existential fresh-message forgery]
For a valid full-rank standard SSD1760 public key, the public affine-completion
algorithm succeeds on each fresh message with probability
\[
 p=\Pr[110\leq\mathrm{Bin}(242,1/2)\leq150]
  =0.9303737782253141.
\]
It uses zero signing queries and has expected public message-query count
$1/p=1.0748368273$.
\end{theorem}

\begin{proof}
Equation~\eqref{eq:completion} is an invertible affine map from the uniform
target $t$ to the pivot vector $s_P$. Hence $s_P$ is uniform in
$\F_2^{242}$, and its Hamming weight has distribution
$\mathrm{Bin}(242,1/2)$. Pivot and free coordinates are disjoint, so
\[
 \wt(s)=363+\wt(s_P).
\]
The inclusive interval in Equation~\eqref{eq:verify} is therefore equivalent
to $110\leq\wt(s_P)\leq150$. Summing the corresponding binomial coefficients
gives
\[
 p=2^{-242}\sum_{j=110}^{150}\binom{242}{j}
  =0.9303737782253141.
\]
The attacker hashes distinct fresh messages and stops at the first accepted
completion. This is a geometric experiment with expectation $1/p$. It never
calls the signing oracle or uses a secret key.
\end{proof}

The method is polynomial-time public linear algebra. It does not solve a
generic exponential low-weight syndrome-decoding problem. The broad central
weight interval makes one fixed affine representative acceptable with
constant probability.

\subsection{Claim boundary}

This theorem is existential: the attacker selects its successful message by
public search. It does not establish a universal signature, success for every
fixed preselected message, or an unconditional property of a particular
concrete hash function. A fixed SHAKE256 run is an executable instance of the
paper-level attack, while the exact probability is a random-oracle statement.

\section{Full-parameter evidence}

\subsection{Primary reconstruction}

The sealed primary Python implementation independently reconstructed the
printed Gen1/Gen2 algebra at complete SSD1760 dimensions. Its valid structured
key had rank 242 for both public matrices. An honest weight-499 signature
passed the ordinary verifier. The public attack fixed 363 free ones, obtained
pivot weight 116 on its first counter message, and produced an accepted
weight-479 response with zero signing queries.

The ordinary verifier rejected the response after changing the message,
changing one response bit while retaining an in-range weight, altering the
public key, and replacing the response by zero. A separately sealed control
reconstructed the same public key and rejected weights 0 and 1760. The full
primary process took 4.87 seconds wall time, 19,968 KiB peak resident memory,
and one worker. The primary result has SHA-256
\begin{center}
\small\ttfamily
87a78f23b8189f27f2339f56a871b4e3871c5128a6f8551953eddaf49b22089d.
\end{center}

\subsection{Hostile adversarial reconstruction}

The clean adversarial implementation was written independently in Magma from
the pinned paper and modeled one fresh uniform 1760-bit oracle response for
each distinct message. It used the paper's consecutive odd signing exponents
$[11,13,15,17]$, obtained a valid rank-242 key on its first Gen1 attempt, and
accepted an honest weight-501 signature.

The public solve tracked the complete row transform, verified the pivot
identity, and agreed with a separate inverse of the original pivot-column
submatrix. Its first candidate had free, pivot, and total weights
$363$, $110$, and $473$ and passed Algorithm~4. Changed-message, in-range
one-bit, and independent structured wrong-key controls rejected.

The adversarial run also constructed responses that satisfied the public
equation at each neighboring endpoint. Table~\ref{tab:endpoints} isolates the
inclusive weight check.

\begin{table}[ht]
\centering
\small
\begin{tabular}{@{}rrrr@{}}
\toprule
Response weight & Public equation & Expected & Ordinary verifier\\
\midrule
472 & holds & reject & reject\\
473 & holds & accept & accept\\
513 & holds & accept & accept\\
514 & holds & reject & reject\\
\bottomrule
\end{tabular}
\caption{Equation-valid boundary controls for Algorithm~4.}
\label{tab:endpoints}
\end{table}

The clean source was sealed before execution at SHA-256
\begin{center}
\small\ttfamily
145d28c510234e76459598b917dc4c55aade5156a68038f019d1196685b431b1.
\end{center}
The run exited zero with empty stderr and no failed assertion. Its attack
stage took 0.19 CPU seconds. The complete process took 1.29 seconds wall time
and 23,552 KiB peak resident memory on one worker. An earlier provisional run
that emitted an interpreter error is explicitly superseded and contributes no
fact or measurement here.

\subsection{Fresh-context independent reproduction}

A separately tasked reproducer worked from the pinned paper and attack
description without reading the primary or adversarial implementations or
results. It independently reconstructed the full standard SSD1760 key,
signing, verification, and public-completion algorithms in Python. Five fresh
full-row-rank structured keys were generated from operating-system entropy,
and all five yielded accepted fresh-message forgeries under the unchanged
ordinary verifier.

The central public-hash query counts were $1,2,1,1,1$, six in total, with zero
signing queries. Forgery weights were $482,489,489,478,478$. Attack times from
immediately before public row reduction were 0.052984--0.053914 seconds per
key and 0.267922 seconds total. The whole five-key process took 14.76 seconds
wall time, 10.66 seconds user CPU, 0.07 seconds system CPU, and 24,576 KiB peak
resident memory on one worker.

For every key, a separate honest signature accepted; changed-message,
weight-preserving changed-signature, and independent wrong-key controls
rejected; and equation-valid weights 472, 473, 513, and 514 respectively
rejected, accepted, accepted, and rejected. A separately sealed validator then
rechecked only the retained public keys, messages, signatures, hashes,
equations, weights, freshness, and query counts and reported a pass. The sealed
independent results have SHA-256
\begin{center}
\small\ttfamily
9b9843dfca266117d3d0a098953816feabb9ce1b19e225cadc79ded93989a878.
\end{center}
The 30-entry experiment manifest has SHA-256
\begin{center}
\small\ttfamily
2608c672de46fd87302c8a8b6d5f3c08b3fb1ffbde39be599e4c8a7facf1abb6.
\end{center}

\subsection{Public package reproducer}

The polished Python reproducer uses only the standard library. It builds two
full structured keys, checks honest acceptance, executes the public attack,
and repeats the dual RREF/pivot-inverse cross-check. It uses SHAKE256 on
explicit distinct message encodings and reports public message and signing
query counts separately. Its controls include a changed message, an in-range
response mutation, an independent wrong key, zero response, and four
equation-valid boundary responses.

The package additionally verifies the separate Gen1 rank defect described
below. Ordinary commands draw fresh operating-system seeds. Disclosed fixed
seeds produce deterministic reference JSON. Timing values are excluded from
those committed JSON files; external measurements must identify either the
attack-stage or complete-process boundary.

\section{Specification issues}

These issues are relevant to interpretation but are not the cause of the
forgery.

\paragraph{Gen1 rank defect.}
The paper permits $\alpha_0=r=15$. With $Q=I+M^r$, the Gen1 middle factor then
becomes
\[
 I+M^{\alpha_0}+R_0Q=Q+R_0Q=(I+R_0)Q.
\]
Its rank is at most $\operatorname{rank}(Q)=198$, below the 242 independent
rows Gen2 requests. Left multiplication cannot increase the rank. Both sealed
attack runs used valid full-rank keys with other $\alpha_0$ values; the public
forgery relies only on those valid keys.

\paragraph{Signing-threshold ambiguity.}
Algorithm~3 prints a while-condition containing both
$\wt(u')<\omega-\mu$ and $\wt(u')<d-\mu$, while the surrounding prose
discusses the $\omega-\mu$ lower bound. Because $d>\omega$, the literal
condition imposes the higher value 491. The sealed honest weights 499 and 501
satisfy both readings, and Algorithm~4 unambiguously accepts from 473 through
513.

\paragraph{Missing concrete interface.}
The paper defines a 1760-bit random oracle but no concrete hash, message
encoding, signature serialization, or author implementation. The attack is
presented at the ordinary mathematical interface. The executable SHAKE256
adapter is explicit and is not described as an author wire format.

\paragraph{SSD2744 inconsistency.}
Table~I prints $k'=377$ for SSD2744, while Tables~III and VI print $k'=337$.
This report makes no claim for SSD2744.

\section{Timer semantics}

The reliable adversarial attack timer starts immediately before public row
reduction and stops after the first candidate passes ordinary verification.
It reports 0.19 \emph{CPU seconds}. The 1.29-second \emph{wall time} covers a
larger process: first structured key generation, honest conformance, the
attack, a second structured wrong key, and endpoint controls. Its peak memory
was 23,552 KiB.

The primary JSON field named \texttt{seconds\_including\_rref} is mislabeled.
Its timer begins after row reduction and the recorded 0.000544 seconds covers
only candidate-message completion. We do not cite it as an attack-stage or
complete-process runtime. CPU and wall measurements, and their different
scopes, are kept separate throughout this package.

\section{Attack history and repair direction}

The mandatory final exact-title, identifier, author, mechanism, citation,
repository, and public-discussion recheck through 14 October 2026 found the
proposal, bibliographic and resource pages, an ordinary downstream citation,
the distinct 2023/1956 FDSD successor~\cite{fdsd}, and unrelated
syndrome-decoding work. It located no previous public exact-target key
recovery, equivalent signer or witness recovery, practical forgery,
substantive parameter compromise, or application of this affine-completion
mechanism to SSD. This moderate-high-confidence result is a bounded public
search and cannot exclude private, deleted, unpublished, or unindexed work.
We make no priority claim and no novelty claim for the underlying linear
algebra.

The ordinary verifier needs a relation whose public affine solution space does
not intersect its accepted response distribution with constant probability.
A repair must reconsider the public dimensions, response constraint, and
security reduction together, then publish a complete hash and encoding
profile and repeat correctness and full-parameter analysis. Merely rejecting
one fixed free assignment has not been shown sufficient. This report does not
validate a drop-in patch.

\section{Reproducibility and publication status}

The reviewed private package contains a landing report, machine-readable
metadata, an exact affected-row record, full source for the standard-library
reproducer, deterministic reference outputs, this \LaTeX{} source and PDF,
and a SHA-256 manifest. The public repository is not modified by this package.

Primary reproduction, hostile adversarial review, fresh-context independent
reproduction on five new full-rank SSD1760 keys, final exact-target history
review, and independent Publication Reviewer review all passed. The package is
classified as \texttt{VERIFIED\_ATTACK} and is ready for exact import; this
review did not modify or publish the public repository.

\section*{Acknowledgements}
The computations were performed on the Norwegian Research and Education Cloud (NREC), using resources provided by the University of Bergen and the University of Oslo.

\begin{thebibliography}{9}

\bibitem{ssd}
Abdelhaliem Babiker.
\newblock Digital Signature from Syndrome Decoding Problem.
\newblock IACR Cryptology ePrint Archive, Report 2022/1698, 2022.
\newblock \url{https://eprint.iacr.org/2022/1698}.

\bibitem{gmr}
Shafi Goldwasser, Silvio Micali, and Ronald L. Rivest.
\newblock A Digital Signature Scheme Secure against Adaptive Chosen-Message
Attacks.
\newblock \emph{SIAM Journal on Computing}, 17(2):281--308, 1988.

\bibitem{fdsd}
Abdelhaliem Babiker.
\newblock A Signature Scheme from Full-Distance Syndrome Decoding.
\newblock IACR Cryptology ePrint Archive, Report 2023/1956, 2023.
\newblock \url{https://eprint.iacr.org/2023/1956}.

\end{thebibliography}

\end{document}

